\chapter{Tympanoplasty}

\section*{What Is This Surgery?}
Tympanoplasty is a surgical procedure performed by an otolaryngologist to repair defects of the tympanic membrane (eardrum) and/or reconstruct the ossicular chain within the middle ear. This intervention is primarily aimed at restoring the integrity of the sound-conducting mechanism, thereby improving hearing and preventing recurrent middle ear infections. It addresses a range of pathologies, including chronic perforations resulting from infection or trauma, and damage to the ossicles caused by disease or injury. By restoring the anatomical and functional integrity of the middle ear, tympanoplasty plays a crucial role in enhancing a patient's auditory function and overall quality of life, while also safeguarding the delicate structures of the inner ear from external pathogens. The procedure can be performed using various techniques, depending on the specific pathology and the surgeon's preference, and it is often accompanied by careful preoperative assessment and postoperative management to optimize outcomes.

\section*{Alternative Names}
\begin{itemize}
  \item Myringoplasty (specifically for tympanic membrane repair)
  \item Ossiculoplasty (specifically for ossicular chain reconstruction)
  \item Middle Ear Reconstruction
  \item Eardrum Grafting
  \item Tympanomastoidectomy (if combined with mastoidectomy)
\end{itemize}

\section*{Relevant Surgical Anatomy}
The surgical anatomy relevant to tympanoplasty includes the tympanic membrane, the ossicular chain, and the surrounding structures of the middle ear. The tympanic membrane separates the external auditory canal from the middle ear and is composed of three layers: the outer epithelial layer, the fibrous middle layer, and the inner mucosal layer. The ossicular chain consists of three small bones: the malleus, incus, and stapes, which transmit sound vibrations from the tympanic membrane to the oval window of the cochlea.

The arterial supply to the tympanic membrane is primarily from the deep auricular artery, a branch of the maxillary artery, and the anterior tympanic artery, a branch of the internal maxillary artery. Venous drainage occurs via the pterygoid plexus and the external jugular vein. The auriculotemporal nerve, a branch of the mandibular nerve (V3), provides sensory innervation to the external ear and the tympanic membrane, while the chorda tympani, a branch of the facial nerve (VII), traverses the middle ear and carries taste fibers.

Lymphatic drainage from the tympanic membrane is primarily to the parotid and cervical lymph nodes. Important surgical landmarks include McBurney's point, which is relevant in other abdominal surgeries, and Calot's triangle, which is significant in gallbladder surgeries but not directly related to tympanoplasty. Understanding these anatomical relationships is crucial for minimizing complications during surgery.

\section*{Surgical Principle \& Rationale}
The fundamental principle of tympanoplasty is the restoration of the middle ear's sound-conducting mechanism, which involves two primary components: repair of the tympanic membrane (myringoplasty) and reconstruction of the ossicular chain (ossiculoplasty). For myringoplasty, a graft, typically autologous tissue such as temporalis fascia, tragal perichondrium, or cartilage, is meticulously harvested and placed to close a perforation in the tympanic membrane. This graft serves as a scaffold, promoting the regrowth of epithelial and fibrous layers to create an intact, functional eardrum capable of vibrating effectively in response to sound waves.

In cases where the ossicular chain is damaged, ossiculoplasty aims to re-establish a continuous and efficient pathway for sound transmission from the tympanic membrane to the oval window of the inner ear. This may involve repairing or replacing eroded, dislocated, or fixed ossicles. Various techniques are employed, including repositioning of existing ossicles, interposition of autologous bone or cartilage, or the use of prosthetic devices such as partial ossicular replacement prostheses (PORP) or total ossicular replacement prostheses (TORP). The ultimate goal is to optimize the mechanical advantage of the ossicular chain, ensuring efficient transfer of acoustic energy to the cochlea.

Key technical goals of tympanoplasty include:

\begin{itemize}
  \item Achieving an intact, mobile, and vibrating tympanic membrane.
  \item Restoring the mechanical advantage and continuity of the ossicular chain.
  \item Creating a well-aerated and infection-free middle ear space.
  \item Preventing recurrent middle ear infections and cholesteatoma formation.
\end{itemize}

\section*{History \& Surgical Evolution}
Early attempts at repairing the tympanic membrane date back to the 17th century, with various materials like paper, rubber, and even pig bladder being explored, albeit with limited success. The modern era of tympanoplasty began in the mid-20th century, driven by advancements in microscopy and surgical techniques.

Pioneering work by German otologists Horst Wullstein in 1952 and Fritz Z\o{}llner in 1953 revolutionized the field. Wullstein introduced a classification system for tympanoplasty based on the extent of middle ear damage and the type of reconstruction required, while Z\o{}llner focused on the importance of an aerated middle ear space. Their contributions established tympanic membrane grafting as a standardized otologic operation. Over the subsequent decades, techniques evolved significantly, moving from skin grafts to more biologically compatible materials like temporalis fascia and cartilage. Further refinements in microscopic surgery, instrumentation, and the development of various ossicular prostheses have continued to improve outcomes and expand the indications for tympanoplasty.

\section*{Clinical Indications}
Tympanoplasty is indicated in various clinical scenarios, including:

\begin{itemize}
  \item To repair a chronic tympanic membrane perforation that has failed to heal spontaneously, leading to recurrent infections or hearing loss.
  \item To address conductive hearing loss caused by ossicular chain discontinuity, erosion, or fixation due to chronic otitis media or trauma.
  \item To manage chronic otitis media with effusion that is refractory to medical management and ventilation tube placement, especially if associated with structural damage.
  \item To remove and reconstruct the middle ear after excision of a cholesteatoma, an abnormal skin growth that can destroy middle ear structures.
  \item To prevent recurrent middle ear infections (otitis media) and associated complications such as mastoiditis, labyrinthitis, or intracranial spread of infection.
  \item To protect the delicate inner ear from external environmental factors like water, bacteria, and foreign bodies through an intact eardrum.
  \item To repair large or persistent traumatic perforations of the tympanic membrane that do not close with conservative management.
  \item To correct tympanic membrane atelectasis or deep retraction pockets that threaten ossicular erosion or the development of cholesteatoma.
\end{itemize}

\section*{Clinical Positioning}
Tympanoplasty is often considered a primary surgical treatment for chronic, stable conditions affecting the tympanic membrane and ossicular chain. For instance, a long-standing, non-healing tympanic membrane perforation in an otherwise healthy ear is a direct indication for primary tympanoplasty. Similarly, conductive hearing loss due to ossicular discontinuity identified after resolution of acute infection may be addressed primarily. 

However, tympanoplasty can also serve as a secondary or staged procedure. This is particularly true in complex cases such as cholesteatoma, where the initial surgery might focus on disease eradication (e.g., mastoidectomy), with reconstruction of the middle ear performed in a subsequent, secondary stage. It may also be considered secondary after failed medical management of chronic otitis media, or as a revision surgery following a previously unsuccessful tympanoplasty.

\section*{Surgical Technique \& Procedure}
Tympanoplasty is typically performed under general anesthesia, though local anesthesia with sedation may be used in select cases. The patient is positioned supine with the head turned to the side, allowing optimal access to the affected ear.

The surgical approach depends on the size and location of the perforation, the extent of middle ear pathology, and surgeon preference. Common approaches include the transcanal (endaural) approach, where instruments are passed through the ear canal, or the postauricular (retroauricular) approach, which involves an incision behind the ear to elevate the skin and muscle, providing a wider view. An endaural incision, made within the external auditory canal, is another variation. Once access is gained, a tympanomeatal flap (skin and fibrous tissue of the ear canal) is carefully elevated to expose the middle ear space and the perforation.

The key operative steps generally include:

\begin{itemize}
  \item **Incision and Flap Elevation:** Creating an incision (transcanal, endaural, or postauricular) and elevating the tympanomeatal flap to expose the middle ear.
  \item **Perforation Edge Preparation:** De-epithelializing the edges of the tympanic membrane perforation to promote graft adherence and healing.
  \item **Graft Harvesting:** Obtaining autologous graft material, most commonly temporalis fascia (from above the ear), tragal perichondrium (from the cartilage in front of the ear canal), or cartilage (from the tragus or concha).
  \item **Graft Placement:** The graft is meticulously placed to cover the perforation. Common techniques include the underlay technique (graft placed medial to the remnant tympanic membrane), overlay technique (graft placed lateral to the remnant), or inlay technique (graft inserted into a prepared groove).
  \item **Ossicular Chain Assessment and Reconstruction:** The ossicular chain (malleus, incus, stapes) is inspected for continuity, erosion, or fixation. If defects are present, reconstruction is performed using autologous ossicles, cartilage, or prosthetic devices (PORP or TORP) to restore sound transmission.
  \item **Flap Repositioning:** The tympanomeatal flap is carefully repositioned over the graft.
  \item **Middle Ear Packing:** Absorbable packing materials, such as Gelfoam or absorbable sponges, are placed in the middle ear space to support the graft and maintain its position during initial healing.
  \item **Closure:** The external incision (if postauricular) is closed with sutures, and the ear canal may be packed with antibiotic-impregnated material.
\end{itemize}

A small drain may be placed temporarily if a postauricular incision was used to prevent hematoma formation.

\section*{Preoperative Workup \& Preparation}
Thorough pre-operative assessment is crucial for a successful tympanoplasty, ensuring the patient is medically optimized and the ear is free from active infection. This preparation typically involves a combination of clinical evaluation, diagnostic tests, and patient education.

\begin{itemize}
  \item **Detailed Otologic Examination:** A comprehensive microscopic examination of the ear is performed to assess the size and location of the perforation, the condition of the middle ear mucosa, and the presence of any cholesteatoma or ossicular damage.
  \item **Audiometry:** A pure-tone audiogram and speech audiometry are essential to quantify the degree and type of hearing loss, providing a baseline for post-operative comparison.
  \item **Imaging Studies:** High-resolution computed tomography (CT) scan of the temporal bone may be indicated for complex cases, such as suspected cholesteatoma, ossicular anomalies, or prior to revision surgery, to delineate middle ear and mastoid anatomy.
  \item **Pre-operative Laboratory Tests:** Standard blood tests (e.g., complete blood count, electrolytes, coagulation profile) and an electrocardiogram (EKG) are performed as per anesthesia protocol to assess overall health and surgical fitness.
  \item **NPO Status:** The patient must remain NPO (nil per os) for 6-8 hours prior to surgery to prevent aspiration during anesthesia.
  \item **Medication Adjustments:** Blood-thinning medications (e.g., aspirin, NSAIDs, warfarin) must be discontinued for a specified period before surgery, as advised by the surgeon and anesthesiologist.
  \item **Antibiotics:** Any active ear infection must be treated and resolved with oral antibiotics or ear drops prior to surgery to minimize the risk of post-operative infection and graft failure. Prophylactic intravenous antibiotics are often administered just before the surgical incision.
  \item **Informed Consent:** A detailed discussion with the surgeon covers the risks, benefits, alternatives, and expected outcomes of the procedure, ensuring the patient provides fully informed consent.
\end{itemize}

\section*{Contraindications \& Precautions}
**Absolute Contraindications:**

\begin{itemize}
  \item Active, uncontrolled purulent middle ear infection or otorrhea (ear discharge) at the time of surgery.
  \item Only hearing ear, especially if the risk of further hearing loss is significant (this is a strong relative contraindication, often treated as absolute).
  \item Severe systemic illness or uncontrolled medical conditions that preclude safe administration of anesthesia or surgical recovery.
  \item Presence of malignancy in the external auditory canal or middle ear requiring alternative management.
  \item Patient's inability or unwillingness to comply with post-operative care instructions.
\end{itemize}

**Relative Contraindications and Precautions:**

\begin{itemize}
  \item Poor Eustachian tube function or a poorly aerated middle ear, which significantly increases the risk of graft failure and recurrent middle ear pathology.
  \item Uncontrolled systemic diseases such as diabetes mellitus, autoimmune disorders, or severe immunosuppression, which can impair wound healing and increase infection risk.
  \item Very large or total tympanic membrane perforations, which may have a higher rate of graft failure.
  \item Previous failed tympanoplasty, which may indicate underlying issues or require a more complex revision approach.
  \item Active upper respiratory tract infection, which can temporarily worsen Eustachian tube function and increase the risk of middle ear infection.
  \item Significant inner ear pathology (e.g., severe sensorineural hearing loss) where the potential for conductive hearing improvement is minimal.
\end{itemize}

\section*{Complications \& Risks}
While tympanoplasty is generally safe, as with any surgical procedure, there are potential risks and complications. These can be broadly categorized into general surgical risks and procedure-specific risks.

**General Surgical Risks:**

\begin{itemize}
  \item **Bleeding:** Minor bleeding is common, but significant hematoma formation requiring intervention is rare.
  \item **Infection:** Wound infection or post-operative middle ear infection can occur, potentially leading to graft failure.
  \item **Anesthesia Complications:** Risks associated with general anesthesia include nausea, vomiting, allergic reactions to medications, respiratory issues, and, rarely, more severe cardiac or neurological events.
\end{itemize}

**Procedure-Specific Risks:**

\begin{itemize}
  \item **Graft Failure or Re-perforation:** The most common complication, where the graft does not heal successfully, resulting in a persistent or recurrent perforation.
  \item **Persistent or Worsened Conductive Hearing Loss:** Despite successful graft take, hearing may not improve as expected, or could even worsen due to scar tissue, graft thickness, or ossicular issues.
  \item **Sensorineural Hearing Loss:** Although rare, damage to the inner ear during surgery can lead to permanent sensorineural hearing loss.
  \item **Tinnitus:** New onset or worsening of pre-existing ringing or buzzing in the ear.
  \item **Dizziness or Vertigo:** Temporary dizziness or imbalance is common post-operatively due to middle ear manipulation, but persistent vertigo is rare.
  \item **Facial Nerve Injury:** Extremely rare but serious complication, potentially leading to temporary or permanent facial weakness or paralysis on the operated side.
  \item **Taste Disturbance (Dysgeusia):** Injury to the chorda tympani nerve, which runs through the middle ear, can cause altered taste sensation on the anterior two-thirds of the tongue, usually temporary.
  \item **Cholesteatoma Recurrence:** If cholesteatoma was present, there is a risk of recurrence, necessitating vigilant follow-up and potentially further surgery.
  \item **Failure of Ossicular Reconstruction:** The reconstructed ossicular chain may not function optimally, requiring revision surgery to improve hearing.
  \item **Allergic Reaction:** Reaction to packing materials, antibiotics, or other substances used during or after surgery.
\end{itemize}

\section*{Postoperative Recovery Timeline}
Immediately after tympanoplasty, the patient is transferred to the post-anesthesia care unit (PACU) for close monitoring as they recover from anesthesia. Pain management is initiated, and anti-nausea medications may be administered. It is common to experience some mild pain, pressure, or a feeling of fullness in the ear, along with temporary dizziness or imbalance. The ear canal will be packed with absorbable material, and an external dressing or cotton ball may be placed. Patients are typically discharged home the same day or the following morning, depending on their recovery and the complexity of the surgery.

During the first 1-2 weeks post-operatively, patients are advised to avoid strenuous activities, heavy lifting, nose blowing, and any activities that increase pressure in the head. Strict water precautions are essential; the ear must be kept dry, meaning no swimming and careful showering with ear protection (e.g., cotton ball with petroleum jelly). Ear drops, often containing antibiotics and steroids, may be prescribed. A follow-up appointment is usually scheduled within 1-2 weeks for initial assessment and, if present, removal of any non-absorbable packing or sutures. Hearing improvement may not be immediately noticeable due to swelling and packing, and it can take several weeks to months for the full auditory benefit to become apparent as the graft heals and the middle ear space re-aerates. Full healing of the tympanic membrane typically takes 2-3 months.

\section*{Surgical Findings \& Pathology}
During tympanoplasty, the surgeon documents the condition of the tympanic membrane, the ossicular chain, and any associated middle ear pathology. The presence of a perforation, its size, and the condition of the surrounding tissue are noted. If cholesteatoma is encountered, its extent and involvement of adjacent structures are carefully assessed. The pathology report on resected tissues, if applicable, represents the nature of the removed tissue, confirming the diagnosis and ruling out malignancy. The successful integration of the graft and the restoration of the tympanic membrane's integrity are key findings that indicate a favorable surgical outcome.

\section*{Expected Postoperative Course}
A normal, successful postoperative course following tympanoplasty is characterized by gradual resolution of pain and discomfort, with patients typically reporting a decrease in ear fullness and pressure within the first week. Hearing improvement may not be immediately noticeable due to swelling and packing, but many patients begin to perceive enhanced auditory function within a few weeks as the graft heals and the middle ear space re-aerates. The tympanic membrane usually heals completely within 2-3 months, at which point audiometric evaluations can be performed to assess functional outcomes. Patients are encouraged to adhere to postoperative care instructions to minimize complications and optimize recovery.

\section*{Postoperative Care \& Follow-up}
Post-operative care is essential for ensuring proper healing and monitoring for complications. Patients are typically scheduled for follow-up visits at 1-2 weeks, 1-3 months, and 6-12 months post-surgery. During these visits, the surgeon will assess the integrity of the tympanic membrane, evaluate hearing improvement, and check for any signs of infection or complications.

\begin{itemize}
  \item **Post-operative Visits:** Initial follow-up typically occurs 1-2 weeks after surgery for removal of any external sutures or non-absorbable ear canal packing. Subsequent visits are scheduled at 1-3 months and 6-12 months to monitor graft healing and middle ear status.
  \item **Audiometry:** A formal audiogram is usually performed 3-6 months post-operatively to objectively assess the final hearing outcome and compare it to pre-operative levels.
  \item **Otoscopic Examinations:** Regular microscopic examinations of the ear are conducted to monitor the integrity of the graft, assess middle ear aeration, and check for any signs of infection, retraction, or cholesteatoma recurrence.
  \item **Activity Resumption:** Activity restrictions are gradually lifted over several weeks, with full return to normal activities, including swimming (with ear protection), typically allowed after 2-3 months, once the graft is fully healed.
  \item **Warning Signs:** Patients are educated on warning signs that necessitate immediate medical attention, such as sudden hearing loss, severe or persistent dizziness, facial weakness, persistent purulent discharge, or fever.
  \item **Revision Surgery:** If graft failure, persistent perforation, or inadequate hearing improvement occurs, revision tympanoplasty may be considered after a suitable healing period, usually 6-12 months.
\end{itemize}

\section*{Epidemiology}
Tympanoplasty is a commonly performed otologic procedure worldwide, with thousands of cases undertaken annually. The incidence varies geographically, often correlating with the prevalence of chronic otitis media, which is higher in developing countries and certain underserved populations. The procedure can be performed across all age groups, from children with chronic perforations or cholesteatoma to adults with long-standing middle ear disease. While there is no significant sex predilection for the underlying conditions requiring tympanoplasty, the specific indications can vary. Chronic tympanic membrane perforation is the most frequent indication, followed by cholesteatoma and ossicular chain defects. The overall mortality associated with tympanoplasty is exceedingly low, primarily linked to general anesthesia risks. Morbidity is mainly related to graft failure (5-15\% recurrence rate), persistent or worsened hearing loss, temporary dizziness, and, rarely, more severe complications like facial nerve injury.

\section*{Outcomes \& Prognosis}
The outcomes of tympanoplasty are generally favorable, with high success rates for anatomical closure of the tympanic membrane and significant improvement in hearing for many patients. Graft take rates for myringoplasty typically range from 85\% to 95\%, depending on factors such as perforation size, presence of infection, and surgical technique. Functional outcomes, specifically hearing improvement, are also positive, with 60\% to 80\% of patients achieving a significant reduction in their air-bone gap. This translates to improved auditory function, reduced frequency of ear infections, and an enhanced quality of life.

The prognosis is generally good, especially for isolated tympanic membrane perforations. Recurrence of perforation is possible, with rates typically between 5\% and 15\%. In cases involving cholesteatoma, the risk of recurrence necessitates long-term follow-up, sometimes including a planned second-stage surgery to check for residual disease. Prognostic factors influencing success include the size of the perforation, the presence and extent of middle ear infection or cholesteatoma, the status of Eustachian tube function, the type of graft material used, and the experience of the surgeon. Patients with good Eustachian tube function and smaller, non-infected perforations generally have the best outcomes.

\section*{References}
\begin{enumerate}
  \item MedlinePlus. Tympanoplasty. U.S. National Library of Medicine. Available at: https://medlineplus.gov/ency/article/002996.htm. Accessed 2024.
  
  \item Flint PW, et al. {\em Cummings Otolaryngology: Head \& Neck Surgery}. 7th ed. Elsevier; 2021. Chapter 151: Tympanoplasty and Ossicular Chain Reconstruction.
  
  \item Paparella MM, et al. {\em Paparella's Otolaryngology}. 3rd ed. Jaypee Brothers Medical Publishers; 2019. Section on Middle Ear Surgery.
  
  \item Rosenfeld RM, et al. Clinical Practice Guideline: Tympanostomy Tubes in Children. {\em Otolaryngology–Head and Neck Surgery}. 2013; 149(1): S1-S35.
\end{enumerate}

\clearpage